%----------------------------------------------------------------------------------------
% KETTINGREGEL EN STANDAARDFUNCTIES. Alle opgaven zijn eigen werk.
%
% Opbouw: Start, leerstof (theorie, werkwijze of uitleg, voorbeeld, korte oefening), en op een
% nieuwe pagina de Opdrachten (door elkaar, stijgende moeilijkheid), Koppel en Test jezelf.
% Labels van dit hoofdstuk beginnen met h3-.
%----------------------------------------------------------------------------------------

\mychapter{3}{De kettingregel}

\lkr{Leerplan B+S'': LPD 26-30 (limieten en afgeleiden), het deel over de kettingregel en de afgeleiden van goniometrische, exponentiële en logaritmische functies. Het bewijs van de kettingregel wordt hier alleen aannemelijk gemaakt. De limiet $\lim_{h \to 0} \frac{\sin h}{h} = 1$ is in Goniometrische functies numeriek verkend; hier wordt ze meetkundig verklaard.}

Met de rekenregels uit hoofdstuk 2 kun je veeltermen en quotiënten afleiden. Maar hoe leid je $(2x + 1)^{10}$ af zonder alles uit te werken? En wat is de afgeleide van $\sin x$, $\e^x$ of $\ln x$? Daarvoor dient dit hoofdstuk. Op het einde heb je een tabel waarmee je elke functie uit deze cursussen kunt afleiden.

\begin{oefeningen}[Start]
\oef{h3-start-1}

\oef{h3-verken}
\end{oefeningen}

%========================================================================================
\section{De kettingregel}

\subsection{Samengestelde functies afleiden}

De functie $h(x) = (2x + 1)^{10}$ is samengesteld: eerst bereken je $u = 2x + 1$, de \textit{binnenste} functie, en daarvan neem je de tiende macht, de \textit{buitenste} functie $u^{10}$. In de Verken zag je dat de afgeleide niet zomaar $10(2x + 1)^9$ is: er komt nog een factor bij, de afgeleide van de binnenste functie.

\begin{theorie}
\begin{stelling}{Kettingregel}
Is $g$ afleidbaar in $x$ en $f$ afleidbaar in $g(x)$, dan is
\[\formule{\bigl(f(g(x))\bigr)' = f'\bigl(g(x)\bigr) \cdot g'(x)}.\]
In de notatie van Leibniz, met $u = g(x)$ en $y = f(u)$: \qquad $\formule{\frac{\mathrm{d}y}{\mathrm{d}x} = \frac{\mathrm{d}y}{\mathrm{d}u} \cdot \frac{\mathrm{d}u}{\mathrm{d}x}}$
\end{stelling}
\end{theorie}

In woorden: leid de buitenste functie af en laat de binnenste ongewijzigd staan, en vermenigvuldig met de afgeleide van de binnenste.

Waarom klopt dat? Verandert $x$ een beetje, dan verandert $u = g(x)$ ongeveer $g'(x)$ keer zoveel, en $y = f(u)$ daardoor nog eens ongeveer $f'(u)$ keer zoveel. De veranderingen worden dus met elkaar vermenigvuldigd. In formules: schrijf het differentiequotiënt van $f(g(x))$ als een product,
\[\frac{f(g(x + h)) - f(g(x))}{h} = \frac{f(g(x + h)) - f(g(x))}{g(x + h) - g(x)} \cdot \frac{g(x + h) - g(x)}{h} .\]
Voor $h \to 0$ gaat de tweede breuk naar $g'(x)$. Omdat $g$ continu is, gaat ook $g(x + h)$ naar $g(x)$, en dan is de eerste breuk een differentiequotiënt van $f$ in $g(x)$, met limiet $f'(g(x))$.

Dat is nog geen volledig bewijs: als $g(x + h) = g(x)$ voor sommige kleine $h$, dan deel je in de eerste breuk door nul. Een volledig bewijs vangt dat op met een technische omweg, die we hier weglaten.

\voorbeeld{}{Bereken de afgeleide van $h(x) = (2x + 1)^{10}$.\label{vb:h3-ketting}}{%
\uitleg{Uitwerken met het binomium kan, maar is veel werk. Herken een samengestelde functie: wat reken je als laatste uit? De tiende macht. Dat is de buitenste functie.}
\begin{stappen}
u = 2x + 1, \quad h = u^{10} & Binnenste en buitenste functie. \\[1ex]
\frac{\mathrm{d}h}{\mathrm{d}u} = 10u^9, \quad \frac{\mathrm{d}u}{\mathrm{d}x} = 2 & Leid elk apart af. \\[1ex]
h'(x) = 10(2x + 1)^9 \cdot 2 = 20(2x + 1)^9 & Kettingregel: vul $u = 2x + 1$ terug in, en vermenigvuldig.
\end{stappen}
Wie de factor $2$ vergeet, krijgt een antwoord dat half zo groot is. Controle in $x = 0$: de grafiek van $h$ is daar twee keer zo steil als die van $u^{10}$ in $u = 1$, omdat $u$ twee keer zo snel verandert als $x$.}

\begin{oefeningen}
\oef{h3-kmacht}
\end{oefeningen}

Bij een wortel is de buitenste functie $\sqrt u$, met afgeleide $\frac{1}{2\sqrt u}$.

\voorbeeld{}{Bereken de afgeleide van $h(x) = \sqrt{x^2 + 1}$.\label{vb:h3-wortel}}{%
\uitleg{Als laatste neem je de wortel: dat is de buitenste functie. De binnenste is $u = x^2 + 1$.}
\begin{stappen}
h'(x) = \frac{1}{2\sqrt{x^2 + 1}} \cdot 2x & De afgeleide van $\sqrt u$, met $u = x^2 + 1$ ingevuld, maal $u' = 2x$. \\[2ex]
= \frac{x}{\sqrt{x^2 + 1}} & Vereenvoudig: de $2$ valt weg.
\end{stappen}}

\begin{oefeningen}
\oef{h3-kwortel}
\end{oefeningen}

\subsection{De notatie van Leibniz}

De notatie $\frac{\mathrm{d}y}{\mathrm{d}x} = \frac{\mathrm{d}y}{\mathrm{d}u} \cdot \frac{\mathrm{d}u}{\mathrm{d}x}$ ziet eruit alsof je $\mathrm{d}u$ wegdeelt. Dat is geen bewijs, want het zijn geen echte breuken, maar het helpt om de regel te onthouden. Ze is vooral handig als de grootheden een eigen naam hebben.

\voorbeeld{}{Neem $y = u^3$ en $u = x^2 + 1$. Bereken $\frac{\mathrm{d}y}{\mathrm{d}x}$, en de waarde ervan in $x = 1$.\label{vb:h3-leibniz}}{%
\uitleg{Bereken de twee afgeleiden apart, vermenigvuldig ze, en schrijf alles in $x$.}
\begin{stappen}
\frac{\mathrm{d}y}{\mathrm{d}u} = 3u^2, \quad \frac{\mathrm{d}u}{\mathrm{d}x} = 2x & $y$ als functie van $u$, en $u$ als functie van $x$. \\[2ex]
\frac{\mathrm{d}y}{\mathrm{d}x} = 3u^2 \cdot 2x = 6x(x^2 + 1)^2 & Kettingregel, met $u = x^2 + 1$ ingevuld. \\[2ex]
x = 1\colon\ u = 2, \ \frac{\mathrm{d}y}{\mathrm{d}x} = 3 \cdot 4 \cdot 2 = 24 & In één punt kun je ook eerst $u$ berekenen.
\end{stappen}}

\subsection{Meer lagen}

Een functie kan uit meer dan twee lagen bestaan. Dan pas je de kettingregel herhaald toe, van buiten naar binnen.

\voorbeeld{}{Bereken de afgeleide van $h(x) = \sqrt{1 + (3x - 1)^2}$.\label{vb:h3-lagen}}{%
\uitleg{Wat reken je als laatste uit? De wortel. Daarbinnen staat $1 + v^2$, en daarbinnen $v = 3x - 1$. Drie lagen, dus drie factoren.}
\begin{stappenbreed}
h'(x) = \frac{1}{2\sqrt{1 + (3x - 1)^2}} \cdot \bigl(1 + (3x - 1)^2\bigr)' & De buitenste laag: de wortel. \\[2ex]
\bigl(1 + (3x - 1)^2\bigr)' = 2(3x - 1) \cdot 3 & De middelste laag: een kwadraat, maal de afgeleide van $3x - 1$. \\[2ex]
h'(x) = \frac{6(3x - 1)}{2\sqrt{1 + (3x - 1)^2}} = \frac{3(3x - 1)}{\sqrt{1 + (3x - 1)^2}} & Samen, en vereenvoudigen.
\end{stappenbreed}}

\begin{werkwijze}
Zoek de buitenste en de binnenste functie: wat reken je als laatste uit? Leid de buitenste af en laat de binnenste ongewijzigd staan. Vermenigvuldig met de afgeleide van de binnenste. Bij meer lagen herhaal je dat, van buiten naar binnen.
\end{werkwijze}

\begin{oefeningen}
\oef{h3-ketting}

\oef{h3-lagen}
\end{oefeningen}

%========================================================================================
\section{Goniometrische functies}

\subsection{Twee limieten}

De afgeleide van $\sin x$ vraagt twee limieten. Werk daarbij altijd in radialen: dan is een hoek $h$ de lengte van de boog die hij op de eenheidscirkel afsnijdt.

\naast{%
Neem $0 < h < \frac{\pi}{2}$. Op de eenheidscirkel liggen $A(1, 0)$, $P(\cos h, \sin h)$ en op de raaklijn in $A$ het punt $T(1, \tan h)$. In de figuur zie je drie stukken die in elkaar passen: de driehoek $OAP$, de cirkelsector $OAP$ en de driehoek $OAT$.

Hun oppervlaktes zijn $\frac12 \cdot 1 \cdot \sin h$, $\frac12 h$ (een sector met hoek $h$ in een cirkel met straal $1$) en $\frac12 \cdot 1 \cdot \tan h$.}{%
\begin{tikzpicture}[scale=2.6,font=\scriptsize]
  \fill[accent!10] (0,0) -- (1,0) -- (1,{tan(40)}) -- cycle;
  \fill[accent!25] (0,0) -- (1,0) arc (0:40:1) -- cycle;
  \fill[accent!45] (0,0) -- (1,0) -- (40:1) -- cycle;
  \draw[grijs] (1,0) arc (0:70:1);
  \draw[->,grijs] (-.05,0) -- (1.2,0);
  \draw[grijs] (1,-.05) -- (1,1);
  \draw[accent,thick] (0,0) -- (1,{tan(40)});
  \draw[accent] (40:1) -- (1,0);
  \node[below left] at (0,0) {$O$};
  \node[below] at (1,0) {$A$};
  \node[above left] at (40:1) {$P$};
  \node[right] at (1,{tan(40)}) {$T$};
  \node at (.3,.1) {$h$};
\end{tikzpicture}}

\begin{theorie}
\begin{stelling}{Limiet van \texorpdfstring{$\frac{\sin h}{h}$}{sin h/h}}
Met $h$ in radialen is \quad $\formule{\lim_{h \to 0} \frac{\sin h}{h} = 1}$.
\end{stelling}
\begin{bewijs}
Neem eerst $0 < h < \frac{\pi}{2}$. De drie oppervlaktes in de figuur passen in elkaar, dus
\bewijsregel{\tfrac12 \sin h < \tfrac12 h < \tfrac12 \tan h .}
Vermenigvuldig met $\frac{2}{\sin h} > 0$: dan is $1 < \frac{h}{\sin h} < \frac{1}{\cos h}$. Neem van alles het omgekeerde (de tekens keren om):
\bewijsregel{\cos h < \frac{\sin h}{h} < 1 .}
Voor $-\frac{\pi}{2} < h < 0$ geldt hetzelfde, want $\frac{\sin(-h)}{-h} = \frac{\sin h}{h}$ en $\cos(-h) = \cos h$. Neem nu $\varepsilon > 0$. Omdat $\cos$ continu is, is $\cos h > 1 - \varepsilon$ voor $h$ dicht genoeg bij $0$, en dan is
\bewijsregel{1 - \varepsilon < \frac{\sin h}{h} < 1 ,}
zodat $\bigl|\frac{\sin h}{h} - 1\bigr| < \varepsilon$. Dat is de definitie van de limiet.
\end{bewijs}
\end{theorie}

Het quotiënt zit ingeklemd tussen $\cos h$ en $1$, en omdat $\cos h$ naar $1$ gaat, kan het nergens anders heen. Meetkundig: voor een kleine hoek zijn de boog $h$ en de verticale $\sin h$ bijna even lang.

\begin{theorie}
\begin{stelling}{Limiet van \texorpdfstring{$\frac{1 - \cos h}{h}$}{(1 - cos h)/h}}
Met $h$ in radialen is \quad $\formule{\lim_{h \to 0} \frac{1 - \cos h}{h} = 0}$.
\end{stelling}
\begin{bewijs}
Vermenigvuldig teller en noemer met $1 + \cos h$, en gebruik $1 - \cos^2 h = \sin^2 h$:
\bewijsregel{\frac{1 - \cos h}{h} = \frac{1 - \cos^2 h}{h\,(1 + \cos h)} = \frac{\sin h}{h} \cdot \frac{\sin h}{1 + \cos h} .}
De eerste factor gaat naar $1$, de tweede naar $\frac{0}{1 + 1} = 0$. Het product gaat naar $1 \cdot 0 = 0$.
\end{bewijs}
\end{theorie}

Het verschil tussen $1$ en $\cos h$ is dus veel kleiner dan $h$ zelf. Met deze twee limieten kunnen we de goniometrische functies afleiden.

\voorbeeld{}{Bereken $\displaystyle\lim_{h \to 0} \dfrac{\sin 3h}{h}$.\label{vb:h3-sinlim}}{%
\uitleg{Invullen geeft $\frac00$. Om de stelling te gebruiken, moet onder $\sin 3h$ ook $3h$ staan. Vermenigvuldig daarom teller en noemer met $3$.}
\begin{stappen}
\frac{\sin 3h}{h} = 3 \cdot \frac{\sin 3h}{3h} & Teller en noemer maal $3$. \\[2ex]
\lim_{h \to 0} 3 \cdot \frac{\sin 3h}{3h} = 3 \cdot 1 = 3 & Stel $u = 3h$: als $h \to 0$, dan $u \to 0$, en $\frac{\sin u}{u} \to 1$.
\end{stappen}}

\begin{oefeningen}
\oef{h3-sinlim}
\end{oefeningen}

\subsection{Afgeleide van de sinus}

\begin{theorie}
\begin{stelling}{Afgeleide van de sinus}
$\formule{(\sin x)' = \cos x}$ \qquad ($x$ in radialen)
\end{stelling}
\begin{bewijs}
Met de somformule $\sin(x + h) = \sin x\cos h + \cos x\sin h$ is
\bewijsregel{\frac{\sin(x + h) - \sin x}{h} = \frac{\sin x\,(\cos h - 1) + \cos x\sin h}{h} = -\sin x \cdot \frac{1 - \cos h}{h} + \cos x \cdot \frac{\sin h}{h} .}
Met de twee limieten gaat dat voor $h \to 0$ naar
\bewijsregel{-\sin x \cdot 0 + \cos x \cdot 1 = \cos x .}
\end{bewijs}
\end{theorie}

Kijk naar de grafieken: waar de sinus het steilst stijgt, in $x = 0$, is de cosinus maximaal. In de toppen van de sinus is de raaklijn horizontaal, en daar is de cosinus nul. Waar de sinus daalt, is de cosinus negatief.

\figuren{\begin{tikzpicture}[x=1cm,y=1cm,font=\scriptsize]
  \draw[black!14,very thin] (-1.6,-1) grid[xstep=1.5708,ystep=1] (6.3,1);
  \draw[->,grijs] (-1.8,0) -- (6.6,0) node[right,text=black] {$x$};
  \draw[->,grijs] (0,-1.3) -- (0,1.5) node[above,text=black] {$y$};
  \foreach \t/\l in {1.5708/\frac{\pi}{2},3.1416/\pi,4.7124/\frac{3\pi}{2},6.2832/2\pi}{\node[below,inner sep=1.5pt] at (\t,0) {$\l$};}
  \node[left,inner sep=1.5pt] at (0,1) {$1$};
  \kromme{-1.57}{6.28}{sin(\x r)}
  \kromme[grijs,dashed]{-1.57}{6.28}{cos(\x r)}
  \draw[grijs,thick] (-.8,-.8) -- (.8,.8);
  \node[accent,above] at (1.57,1) {$\sin x$};
  \node[grijs,below] at (3.14,-1) {$\cos x$};
\end{tikzpicture}}

\voorbeeld{}{Geef de raaklijn aan de grafiek van $f(x) = \sin x$ in de oorsprong.\label{vb:h3-sinraak}}{%
\uitleg{De richtingscoëfficiënt is de afgeleide in $0$.}
\begin{stappen}
f(0) = 0, \quad f'(0) = \cos 0 = 1 & De afgeleide van de sinus is de cosinus. \\[1ex]
y = 0 + 1 \cdot (x - 0) = x & De raaklijn in de figuur.
\end{stappen}
Dicht bij $0$ is dus $\sin x \approx x$, zolang $x$ in radialen staat. Dat is de limiet $\frac{\sin h}{h} \to 1$ in een andere vorm.}

\begin{oefeningen}
\oef{h3-sin}
\end{oefeningen}

\subsection{Afgeleide van de cosinus}

De cosinus kun je net zo afleiden, met de somformule voor $\cos(x + h)$: dat is een oefening hieronder. Met de kettingregel gaat het sneller, want de cosinus is een verschoven sinus.

\begin{theorie}
\begin{stelling}{Afgeleide van de cosinus}
$\formule{(\cos x)' = -\sin x}$ \qquad ($x$ in radialen)
\end{stelling}
\begin{bewijs}
Gebruik $\cos x = \sin\bigl(\frac{\pi}{2} - x\bigr)$ en $\sin x = \cos\bigl(\frac{\pi}{2} - x\bigr)$ (complementaire hoeken). De kettingregel, met binnenste functie $\frac{\pi}{2} - x$ en afgeleide $-1$, geeft
\bewijsregel{(\cos x)' = \cos\bigl(\tfrac{\pi}{2} - x\bigr) \cdot (-1) = -\sin x .}
\end{bewijs}
\end{theorie}

Let op het min-teken: de cosinus daalt tussen $0$ en $\pi$, precies waar de sinus positief is.

\voorbeeld{}{Bereken de afgeleide van $g(x) = \cos 3x$.\label{vb:h3-cos}}{%
\uitleg{De binnenste functie is $3x$, met afgeleide $3$. De buitenste is $\cos u$, met afgeleide $-\sin u$.}
\begin{stappen}
g'(x) = -\sin 3x \cdot 3 = -3\sin 3x & Kettingregel: de factor $3$ komt ervoor.
\end{stappen}
Vergelijk de grafieken: $\cos 3x$ heeft een drie keer kleinere periode dan $\cos x$, en is daardoor ook drie keer zo steil.}

\begin{oefeningen}
\oef{h3-cosbewijs}
\end{oefeningen}

\subsection{Afgeleide van de tangens}

De tangens is een quotiënt, $\tan x = \frac{\sin x}{\cos x}$. Met de quotiëntregel volgt haar afgeleide uit die van de sinus en de cosinus.

\begin{theorie}
\begin{stelling}{Afgeleide van de tangens}
$\formule{(\tan x)' = \frac{1}{\cos^2 x} = 1 + \tan^2 x}$ \qquad (waar $\cos x \neq 0$)
\end{stelling}
\begin{bewijs}
Met de quotiëntregel en de grondformule $\cos^2 x + \sin^2 x = 1$ is
\bewijsregel{\Bigl(\frac{\sin x}{\cos x}\Bigr)' = \frac{\cos x \cdot \cos x - \sin x \cdot (-\sin x)}{\cos^2 x} = \frac{\cos^2 x + \sin^2 x}{\cos^2 x} .}
Dat is $\frac{1}{\cos^2 x}$. Splits je de breuk, dan krijg je $1 + \frac{\sin^2 x}{\cos^2 x} = 1 + \tan^2 x$.
\end{bewijs}
\end{theorie}

De afgeleide is overal positief: de tangens stijgt op elk interval van haar domein, en wordt steeds steiler naar de asymptoten toe.

\voorbeeld{}{Bereken de afgeleide van $f(x) = \tan 2x$, en de richtingscoëfficiënt van de raaklijn in $x = \frac{\pi}{8}$.\label{vb:h3-tan}}{%
\uitleg{Kettingregel, met binnenste functie $2x$.}
\begin{stappen}
f'(x) = \frac{1}{\cos^2 2x} \cdot 2 = \frac{2}{\cos^2 2x} & De afgeleide van $\tan u$, maal $u' = 2$. \\[2ex]
f'\bigl(\tfrac{\pi}{8}\bigr) = \frac{2}{\cos^2 \frac{\pi}{4}} = \frac{2}{\frac12} = 4 & $\cos \frac{\pi}{4} = \frac{\sqrt2}{2}$, met kwadraat $\frac12$.
\end{stappen}}

\begin{oefeningen}
\oef{h3-gonio}
\end{oefeningen}

\subsection{Machten van goniometrische functies}

Bij $\sin^2 x$ staat de macht buiten: $\sin^2 x = (\sin x)^2$. De buitenste functie is dus het kwadraat, en de binnenste de sinus.

\voorbeeld{}{Bereken de afgeleide van $f(x) = \sin^2 x$.\label{vb:h3-gonio}}{%
\uitleg{Buitenste functie $u^2$, binnenste $u = \sin x$.}
\begin{stappen}
f'(x) = 2\sin x \cdot \cos x & De afgeleide van $u^2$ is $2u$, maal $(\sin x)' = \cos x$. \\[1ex]
= \sin 2x & De verdubbelingsformule maakt het korter.
\end{stappen}
Verwar $\sin^2 x$ niet met $\sin x^2 = \sin(x^2)$. Daar is de sinus de buitenste functie, en is de afgeleide $\cos x^2 \cdot 2x$.}

\begin{oefeningen}
\oef{h3-gonio2}
\end{oefeningen}

%========================================================================================
\section{Exponentiële en logaritmische functies}

\subsection{Afgeleide van \texorpdfstring{$\e^x$}{e tot de x}}

Voor elke exponentiële functie $a^x$ (met $a > 0$) kun je het differentiequotiënt vereenvoudigen met een rekenregel voor machten:
\[\frac{a^{x + h} - a^x}{h} = \frac{a^x \cdot a^h - a^x}{h} = a^x \cdot \frac{a^h - 1}{h} .\]
De laatste breuk hangt niet af van $x$: het is het differentiequotiënt van $a^x$ in $0$. De afgeleide van $a^x$ is dus $a^x$ maal de helling van de grafiek in $(0, 1)$. Een rekentoestel geeft voor $h = 0{,}0001$:

\naast{%
\figuren{\renewcommand{\arraystretch}{1.2}%
\begin{tabular}{c|ccc}
$a$ & $2$ & $\e \approx 2{,}718$ & $3$ \\ \hline
$\frac{a^h - 1}{h}$ & $0{,}6932$ & $1{,}0001$ & $1{,}0987$
\end{tabular}}
Het getal $\e$ is precies het grondtal waarvoor die helling $1$ is: de grafiek van $\e^x$ raakt in $(0, 1)$ de rechte $y = x + 1$. In de figuur ligt de grafiek van $\e^x$ tussen die van $2^x$ en $3^x$, en raakt alleen zij die rechte.}{%
\begin{tikzpicture}[x=.8cm,y=.8cm,font=\scriptsize]
  \rooster{-2}{1}{1}{-1}{1}{3}
  \kromme[grijs]{-2}{1.4}{2^(\x)}
  \kromme[grijs]{-2}{1}{3^(\x)}
  \kromme[very thick]{-2}{1.1}{exp(\x)}
  \draw[grijs,dashed] (-2,-1) -- (1.4,2.4);
  \node[grijs,right] at (1.4,2.64) {$2^x$};
  \node[grijs,left] at (.95,3) {$3^x$};
  \vol{(0,1)}
\end{tikzpicture}}

\begin{theorie}
\begin{stelling}{Afgeleide van \texorpdfstring{$\e^x$}{e tot de x}}
$\formule{(\e^x)' = \e^x}$
\end{stelling}
\end{theorie}

Dat $\lim_{h \to 0} \frac{\e^h - 1}{h}$ precies $1$ is, volgt uit de definitie van $\e$ als limiet van $\bigl(1 + \frac1n\bigr)^n$. Dat bewijs is te zwaar voor dit boek; de tabel maakt het aannemelijk. Uit die ene limiet volgt de stelling met de berekening hierboven, voor $a = \e$.

De functie $\e^x$ is dus haar eigen afgeleide: in elk punt is de helling van de grafiek gelijk aan de hoogte. Daarom heet $\e$ het natuurlijke grondtal.

\voorbeeld{}{Bereken de afgeleide van $f(x) = \e^{-x^2}$.\label{vb:h3-explog}}{%
\uitleg{De buitenste functie is $\e^u$, en die is haar eigen afgeleide. De binnenste is $u = -x^2$.}
\begin{stappen}
f'(x) = \e^{-x^2} \cdot (-2x) & Kettingregel: $\e^u$ blijft staan, maal $u' = -2x$. \\[1ex]
= -2x\,\e^{-x^2} & Zet de factor vooraan.
\end{stappen}
Omdat $\e^{-x^2} > 0$, heeft $f'(x)$ het teken van $-2x$: de functie stijgt voor $x < 0$ en daalt voor $x > 0$.}

\begin{oefeningen}
\oef{h3-exp}
\end{oefeningen}

\subsection{Afgeleide van \texorpdfstring{$a^x$}{a tot de x}}

Voor een ander grondtal schrijf je de macht met grondtal $\e$: omdat $a = \e^{\ln a}$, is $a^x = \e^{x \ln a}$.

\begin{theorie}
\begin{stelling}{Afgeleide van \texorpdfstring{$a^x$}{a tot de x}}
Voor $a > 0$: \quad $\formule{(a^x)' = a^x \ln a}$
\end{stelling}
\begin{bewijs}
Schrijf $a^x = \e^{x \ln a}$. De binnenste functie $x \ln a$ heeft afgeleide $\ln a$, een constante. De kettingregel geeft
\bewijsregel{(a^x)' = \e^{x \ln a} \cdot \ln a = a^x \ln a .}
\end{bewijs}
\end{theorie}

De factor $\ln a$ is de helling van $a^x$ in $(0, 1)$. Vergelijk met de tabel: $\ln 2 \approx 0{,}6931$ en $\ln 3 \approx 1{,}0986$. Voor $a = \e$ is $\ln \e = 1$, en krijg je de vorige stelling terug.

\voorbeeld{}{Bereken de afgeleide van $f(x) = 3^{2x + 1}$.\label{vb:h3-ax}}{%
\uitleg{Een exponentiële functie met grondtal $3$, en in de exponent een binnenste functie. Dus de nieuwe regel en de kettingregel.}
\begin{stappen}
f'(x) = 3^{2x + 1} \ln 3 \cdot 2 & De afgeleide van $3^u$ is $3^u \ln 3$, maal $u' = 2$. \\[1ex]
= 2\ln 3 \cdot 3^{2x + 1} & De constante factoren vooraan.
\end{stappen}}

\begin{oefeningen}
\oef{h3-ax}
\end{oefeningen}

\subsection{Afgeleide van \texorpdfstring{$\ln x$}{ln x}}

De natuurlijke logaritme is de inverse van $\e^x$: $\e^{\ln x} = x$ voor $x > 0$. Leid je beide leden af, dan volgt de afgeleide van $\ln x$.

\begin{theorie}
\begin{stelling}{Afgeleide van \texorpdfstring{$\ln x$}{ln x}}
Voor $x > 0$: \quad $\formule{(\ln x)' = \frac{1}{x}}$
\end{stelling}
\begin{bewijs}
Leid beide leden van $\e^{\ln x} = x$ af. Links gebruik je de kettingregel, met binnenste functie $\ln x$:
\bewijsregel{\e^{\ln x} \cdot (\ln x)' = 1 , \qquad\text{dus}\qquad x \cdot (\ln x)' = 1 .}
Deel door $x$: $(\ln x)' = \frac1x$.
\end{bewijs}
\end{theorie}

Dit bewijs gaat ervan uit dat $\ln x$ afleidbaar is. Dat is aannemelijk: de grafiek van $\ln x$ is het spiegelbeeld van die van $\e^x$ om de rechte $y = x$, en heeft dus net als die grafiek overal een raaklijn. Merk op dat de afgeleide van $\ln x$ geen logaritme meer bevat: een eenvoudige macht, $x^{-1}$.

\voorbeeld{}{Bereken de afgeleide van $g(x) = \ln(x^2 + 1)$.\label{vb:h3-ln}}{%
\uitleg{De buitenste functie is $\ln u$, met afgeleide $\frac1u$. De binnenste is $u = x^2 + 1$, en die is altijd positief.}
\begin{stappen}
g'(x) = \frac{1}{x^2 + 1} \cdot 2x = \frac{2x}{x^2 + 1} & Kettingregel.
\end{stappen}
Bij $\ln u$ geeft de kettingregel altijd $\frac{u'}{u}$: de afgeleide van de binnenste functie gedeeld door de binnenste functie.}

\begin{oefeningen}
\oef{h3-ln}
\end{oefeningen}

\subsection{Afgeleide van \texorpdfstring{$\log_a x$}{log x}}

Een logaritme met een ander grondtal schrijf je met de natuurlijke logaritme: $\log_a x = \frac{\ln x}{\ln a}$.

\begin{theorie}
\begin{stelling}{Afgeleide van \texorpdfstring{$\log_a x$}{log x}}
Voor $a > 0$, $a \neq 1$ en $x > 0$: \quad $\formule{(\log_a x)' = \frac{1}{x \ln a}}$
\end{stelling}
\begin{bewijs}
Uit $\log_a x = \frac{1}{\ln a} \cdot \ln x$ volgt met de veelvoudregel
\bewijsregel{(\log_a x)' = \frac{1}{\ln a} \cdot \frac{1}{x} = \frac{1}{x \ln a} .}
\end{bewijs}
\end{theorie}

\voorbeeld{}{Bereken de afgeleide van $f(x) = \log_3(2x - 1)$.\label{vb:h3-log}}{%
\uitleg{De buitenste functie is $\log_3 u$, de binnenste $u = 2x - 1$.}
\begin{stappen}
f'(x) = \frac{1}{(2x - 1)\ln 3} \cdot 2 = \frac{2}{(2x - 1)\ln 3} & De afgeleide van $\log_3 u$ is $\frac{1}{u \ln 3}$, maal $u' = 2$.
\end{stappen}}

\begin{oefeningen}
\oef{h3-explog2}
\end{oefeningen}

\subsection{Machten met een reële exponent}

In hoofdstuk 2 bewezen we de regel voor machten alleen voor natuurlijke en negatieve gehele exponenten, en voor $\frac12$. Met $\e^x$ en $\ln x$ kunnen we ze nu voor elke exponent bewijzen, als $x > 0$.

\begin{theorie}
\begin{stelling}{Afgeleide van een macht, voor \texorpdfstring{$x > 0$}{x > 0}}
Voor elk reëel getal $r$ en $x > 0$: \quad $\formule{(x^r)' = r\,x^{r - 1}}$
\end{stelling}
\begin{bewijs}
Voor $x > 0$ is $x = \e^{\ln x}$, dus $x^r = \e^{r \ln x}$. De kettingregel geeft
\bewijsregel{(x^r)' = \e^{r \ln x} \cdot \frac{r}{x} = x^r \cdot \frac{r}{x} = r\,x^{r - 1} .}
\end{bewijs}
\end{theorie}

Zo is de regel voor machten voor $x > 0$ volledig bewezen, ook voor exponenten zoals $\frac23$ of $\sqrt2$.

\subsection{Combinaties}

Met de rekenregels uit hoofdstuk 2 combineer je de nieuwe afgeleiden.

\voorbeeld{}{Bereken de afgeleide van $f(x) = x^2\e^x$.\label{vb:h3-combinatie}}{%
\uitleg{Een product van twee factoren, dus de productregel. Elke factor kunnen we afleiden.}
\begin{stappen}
f'(x) = 2x \cdot \e^x + x^2 \cdot \e^x & De eerste factor afgeleid maal de tweede, plus de eerste maal de tweede afgeleid. \\[1ex]
= (x^2 + 2x)\,\e^x & Zet $\e^x$ buiten haakjes: dan zie je makkelijker waar $f'$ nul is.
\end{stappen}}

\voorbeeld{}{Bereken de afgeleide van $g(x) = \dfrac{\ln x}{x}$.\label{vb:h3-combinatie2}}{%
\uitleg{Een quotiënt, dus de quotiëntregel, met teller $\ln x$ en noemer $x$.}
\begin{stappen}
g'(x) = \frac{\frac1x \cdot x - \ln x \cdot 1}{x^2} & Quotiëntregel. \\[2ex]
= \frac{1 - \ln x}{x^2} & $\frac1x \cdot x = 1$.
\end{stappen}}

\begin{oefeningen}
\oef{h3-explog}

\oef{h3-combi}
\end{oefeningen}

\subsection{Overzicht}

Met de rekenregels uit hoofdstuk 2 en de afgeleiden uit dit hoofdstuk kun je nu elke functie uit deze cursus afleiden. In de rechterkolom staat telkens de vorm met de kettingregel, met $u$ een functie van $x$.

\needspace{18\baselineskip}
\figuren{\renewcommand{\arraystretch}{2.4}%
\begin{tabular}{|c|c|c|}
\hline
\textbf{functie} & \textbf{afgeleide} & \textbf{met de kettingregel} \\ \hline
$x^r$ & $r\,x^{r-1}$ & $(u^r)' = r\,u^{r-1}\,u'$ \\ \hline
$\sin x$ & $\cos x$ & $(\sin u)' = \cos u \cdot u'$ \\ \hline
$\cos x$ & $-\sin x$ & $(\cos u)' = -\sin u \cdot u'$ \\ \hline
$\tan x$ & $\dfrac{1}{\cos^2 x} = 1 + \tan^2 x$ & $(\tan u)' = \dfrac{u'}{\cos^2 u}$ \\ \hline
$\e^x$ & $\e^x$ & $(\e^u)' = \e^u\,u'$ \\ \hline
$a^x$ & $a^x \ln a$ & $(a^u)' = a^u \ln a \cdot u'$ \\ \hline
$\ln x$ & $\dfrac1x$ & $(\ln u)' = \dfrac{u'}{u}$ \\ \hline
$\log_a x$ & $\dfrac{1}{x \ln a}$ & $(\log_a u)' = \dfrac{u'}{u \ln a}$ \\ \hline
\end{tabular}}

%========================================================================================
\opdrachtenkop

\begin{oefeningen}
\needspace{12\baselineskip}
\oef{h3-o9}

\oef{h3-o2}

\needspace{12\baselineskip}
\oef{h3-o10}

\needspace{8\baselineskip}
\oef{h3-o12}

\oef{h3-o1}

\oef{h3-o13}

\needspace{12\baselineskip}
\oef{h3-o11}

\needspace{14\baselineskip}
\oef{h3-o14}

\oef{h3-o3}

\needspace{6\baselineskip}
\oef{h3-o15}

\needspace{14\baselineskip}
\oef{h3-o16}

\oef{h3-o5}

\oef{h3-o18}

\oef{h3-o4}

\oef{h3-o21}

\needspace{10\baselineskip}
\oef{h3-o17}

\oef{h3-o22}

\oef{h3-o24}

\needspace{10\baselineskip}
\oef{h3-o19}

\oef{h3-o20}

\oef{h3-o6}

\oef{h3-o23}

\oef{h3-o7}

\oef{h3-o25}

\oef{h3-o8}

\oef{h3-o27}

\oef{h3-o26}

\oef{h3-o28}
\end{oefeningen}

%----------------------------------------------------------------------------------------
% KOPPEL: verbanden met andere leerstof
%----------------------------------------------------------------------------------------
\needspace{20\baselineskip}
\opdrachtsectie{Koppel}

\lkr{Integralen, Exponentiële en logaritmische functies en Goniometrische functies: kijk na wat al gezien is. Het pakket legt de volgorde niet vast.}

\begin{oefeningen}
\oef{h3-koppel-1}

\oef{h3-koppel-2}

\oef{h3-koppel-3}
\end{oefeningen}

%----------------------------------------------------------------------------------------
\opdrachtsectie{Test jezelf}

\begin{oefeningen}
\oef{h3-t1}

\oef{h3-t2}

\oef{h3-t3}

\oef{h3-t7}

\oef{h3-t4}

\oef{h3-t5}

\oef{h3-t8}

\oef{h3-t6}
\end{oefeningen}

\herhaalschema{%
  \fv{oef:h3-t1}{vb:h3-ketting}{oef:h3-kmacht}%
  \fv{oef:h3-t2}{vb:h3-gonio}{oef:h3-gonio}%
  \fv{oef:h3-t3}{vb:h3-explog}{oef:h3-o11}%
  \fv{oef:h3-t7}{vb:h3-ax}{oef:h3-explog2}%
  \fv{oef:h3-t4}{vb:h3-combinatie}{oef:h3-combi}%
  \fv{oef:h3-t5}{vb:h3-wortel}{oef:h3-o15}%
  \fzv{oef:h3-t8}{oef:h3-o17}%
  \fzv{oef:h3-t6}{oef:h3-o5}}

\ifzelfstudie
\section*{Maak je eigen samenvatting}
Sluit dit boek en schrijf in je schrift: de kettingregel in woorden en in formule, de afgeleiden van $\sin$, $\cos$, $\tan$, $\e^x$, $a^x$, $\ln x$ en $\log_a x$, en waarom $\frac{\sin h}{h} \to 1$. Maak er een tabel van, samen met de regels uit hoofdstuk 2. Vergelijk daarna met de kaders.
\fi
